% Part: propositional-logic

\documentclass[../../include/open-logic-part]{subfiles}

\begin{document}

\olpart{pl}{ప్రవచన తర్కం}

\begin{editorial}
  ఈ భాగంలో శాస్త్రీయ ప్రవచన తర్కానికి సంబంధించిన విషయం ఉంది. మొదటి
  అధ్యాయం సాపేక్షంగా ప్రాథమికమైనది; అది నిర్వచనాలు, ఫలితాలను మాత్రమే
  జాబితా చేస్తుంది. అనేక నిరూపణలను పూర్తిగా ఇవ్వకుండా సాధనలుగా
  వదిలాం. నిరూపణ వ్యవస్థలు, సంపూర్ణతా సిద్ధాంతానికి సంబంధించిన
  విషయాన్ని ప్రథమక్రమ తర్కం భాగం నుంచి, “FOL” ట్యాగును falseగా
  పెట్టి చేర్చాం. దీనివల్ల విధేయాలు, పదాలు, పరిమాణీకరణలకు సంబంధించిన
  ప్రతిదీ తొలగి, !!{structure}s~$\Struct{M}$ గురించిన చర్చ స్థానంలో
  !!{valuation}s~$\pAssign{v}$ గురించిన చర్చ వస్తుంది.

  ఈ భాగాన్ని మరిన్ని వివరాలతో విస్తరించి, సత్యమూల్య-ప్రమేయ
  సంపూర్ణత వంటి అదనపు అంశాలు, ఫలితాలను చేర్చాలని యోచించాం.
\end{editorial}

\olimport[syntax-and-semantics]{syntax-and-semantics}

\tagfalse{FOL}

\olimport[../first-order-logic/proof-systems]{proof-systems}

\iftag{prfSC}{%
  \olimport[../first-order-logic/sequent-calculus]{sequent-calculus}
}{}

\iftag{prfND}{%
  \olimport[../first-order-logic/natural-deduction]{natural-deduction}
}{}

\iftag{prfTab}{%
  \olimport[../first-order-logic/tableaux]{tableaux}
}{}

\iftag{prfAX}{%
  \olimport[../first-order-logic/axiomatic-deduction]{axiomatic-deduction}
}{}

\olimport[../first-order-logic/completeness]{completeness}

\tagtrue{FOL}

\OLEndPartHook

\end{document}
